
% =============================================================
\section{Applications of Compton Scattering}
% =============================================================

The Compton effect is not only of fundamental physical importance,
but also has a wide range of practical applications in research,
medicine, astronomy, and technology. 
The following sections present some of its most important uses.

% -------------------------------------------------------------
\subsection{Compton Scattering Spectroscopy}
% -------------------------------------------------------------
One of the most direct applications is \textbf{Compton scattering spectroscopy}.
Here, the scattering of high-energy photons by electrons is used to obtain
information about the momentum distribution of electrons in a material.
This method provides insight into the electron density and bonding structure
in solids, liquids, or molecules.

\[
\Delta E = E - E' = E \left(1 - \frac{1}{1 + \frac{E}{m_e c^2}(1 - \cos \theta)}\right)
\]

The measured energy shift serves as the basis for precise
material characterization.

% -------------------------------------------------------------
\subsection{Astrophysics and Cosmology}
% -------------------------------------------------------------
In astrophysics, Compton scattering plays a central role.
In so-called \emph{inverse Compton scattering}, fast electrons transfer
energy to low-energy photons, which are thereby boosted into the
X-ray or gamma-ray range.

These processes explain many phenomena observed in
\textbf{quasars}, \textbf{pulsars}, \textbf{supernova remnants}, and
\textbf{active galactic nuclei}.
A particularly important special case is the \textbf{Sunyaev–Zel’dovich effect},
in which photons of the cosmic microwave background gain energy through
scattering by hot electrons in galaxy clusters.
This allows large-scale structures of the universe to be studied.

% -------------------------------------------------------------
\subsection{Medical Imaging}
% -------------------------------------------------------------
In medical diagnostics, Compton scattering appears both as a
disturbance and as a useful signal.
It is one of the main causes of image blurring in X-ray and CT imaging,
which is why modern systems apply correction algorithms.

Moreover, it is deliberately used in the \textbf{Compton camera}
to reconstruct the spatial distribution of radioactive sources —
a technique applied in nuclear medicine and radiation protection.

% -------------------------------------------------------------
\subsection{Detector Technology and Radiation Protection}
% -------------------------------------------------------------
In the energy range between \SI{100}{\kilo\electronvolt} and \SI{10}{\mega\electronvolt},
Compton scattering is the dominant interaction mechanism
of photons with matter.
It determines the shape of the so-called \emph{Compton continuum}
in gamma spectra and is crucial for the calibration of scintillation
and semiconductor detectors.

It also plays an important role in determining the direction of photon sources
and in material studies using scattering methods.

% -------------------------------------------------------------
\subsection{Material Analysis and Non-Destructive Testing}
% -------------------------------------------------------------
By measuring the energy distribution of scattered photons,
elemental compositions and electron densities can be determined.
This method is particularly suited for light elements,
where the photoelectric effect and pair production are less significant.
Applications are found in materials testing, archaeometry,
and in the analysis of planetary surfaces by space probes.

% -------------------------------------------------------------
\subsection{Fundamentals of Quantum Physics}
% -------------------------------------------------------------
The Compton effect remains to this day a key result of quantum physics:
It directly demonstrates that photons possess energy and momentum
and can interact with particles.
Many teaching and demonstration experiments use Compton scattering
to illustrate the principles of energy and momentum conservation
at the quantum level.

% -------------------------------------------------------------
\begin{PhysicsBox}[Summary]
	Compton scattering connects fundamental quantum physics with practical
	applications — from spectroscopy and astrophysics to medicine.
	It shows that the same physical laws underlying light quanta
	become visible in macroscopic technologies.
\end{PhysicsBox}
